\chapter{Where Latency Comes From}\label{ch:ll:where-latency-comes-from}

A study of a large equity exchange's message data found that races to trade on a stale quote, or to cancel it
before someone else trades on it, happen about once a minute in each of the largest stocks, account for about a
fifth of the day's volume, and are usually decided in five to ten millionths of a second (Aquilina, Budish and
O'Neill, 2022). In those races a microsecond is a unit of competition. This book is about the software on either
side of it: where each microsecond of a trading system goes, how to measure it, and how to take it back. This first
chapter fixes the vocabulary: the path a market event takes through a trading system, why its duration is a
distribution rather than a number, how to write a budget for it, and which strategies pay for which part of it.

\section{The tick-to-trade path}

\begin{definition}[Latency]\label{def:ll:where-latency-comes-from:latency}
The \emph{latency}\index{latency} of a system between two events is the time from the first event (a packet
arriving, a function being called) to the second (an order leaving, the function returning), measured for one
occurrence. It is a random variable: each message has its own.
\end{definition}

One Quant Book~7, chapter~18, used two latencies of a whole trading firm, the market-data \omterm{def:ll:where-latency-comes-from:latency}{latency} and the
order-entry \omterm{def:ll:where-latency-comes-from:latency}{latency}, to model what a strategy sees and when its orders act. This book opens them up.

\begin{definition}[Tick-to-trade and wire-to-wire latency]\label{def:ll:where-latency-comes-from:t2t}
The \emph{tick-to-trade latency}\index{tick-to-trade latency} of a trading system is the time from the arrival of
the market-data message that triggers a decision to the departure of the order that the decision sends. It is a
\emph{wire-to-wire latency}\index{wire-to-wire latency} when both instants are taken on the network cable (the last
bit of the inbound packet, the first bit of the outbound one), by a device outside the server; it is a software
tick-to-trade \omterm{def:ll:where-latency-comes-from:latency}{latency} when they are taken inside the program (the packet handed to the application, the order handed
to the network stack).
\end{definition}

The difference matters. The wire-to-wire figure includes the network card, the operating system's receive path and
the card's transmit path; the software figure, which a program can measure on itself with the processor's clock
(chapter~5), leaves them out. A vendor who quotes a \omterm{def:ll:where-latency-comes-from:latency}{latency} without saying between which two instants has not
quoted a \omterm{def:ll:where-latency-comes-from:latency}{latency}.

\begin{omfigure}
\begin{tikzpicture}[font=\footnotesize,
  st/.style={draw=omDef,fill=omDef!8,rounded corners=2pt,minimum height=0.9cm,text width=1.3cm,align=center},
  hw/.style={draw=gray,fill=gray!10,rounded corners=2pt,minimum height=0.9cm,text width=1.1cm,align=center}]
\node[hw] (nic1) at (0,0) {network card};
\node[hw] (os) at (1.5,0) {kernel or bypass};
\node[st] (dec) at (3.2,0) {decode, arbitrate};
\node[st] (book) at (4.9,0) {book};
\node[st] (strat) at (6.6,0) {strategy};
\node[st] (risk) at (8.3,0) {risk checks};
\node[st] (gw) at (10.0,0) {encode, send};
\node[hw] (nic2) at (11.6,0) {network card};
\foreach \a/\b in {nic1/os,os/dec,dec/book,book/strat,strat/risk,risk/gw,gw/nic2} \draw[-Stealth,thick] (\a) -- (\b);
\draw[-Stealth,thick] (-1.1,0) node[left,align=right,font=\footnotesize]{packet\\ in} -- (nic1);
\draw[-Stealth,thick] (nic2) -- (12.7,0) node[right,align=left,font=\footnotesize]{order\\ out};
\draw[decorate,decoration={brace,amplitude=5pt}] (2.5,0.6) -- (10.7,0.6) node[midway,above=6pt]{software tick-to-trade (the hot path)};
\draw[decorate,decoration={brace,amplitude=5pt,mirror}] (-0.6,-0.6) -- (12.2,-0.6) node[midway,below=6pt]{wire-to-wire, timestamped by a capture device on the cable};
\end{tikzpicture}

\omcaption{The tick-to-trade path of a trading server. The software stages (blue) are this book's subject; the
network card and the kernel's receive and transmit paths (grey) are treated from the program's side in chapters~4
and~13, and in hardware in One Quant Book~14.}\label{fig:ll:where-latency-comes-from:path}
\end{omfigure}

\begin{definition}[Hot path]\label{def:ll:where-latency-comes-from:hot}
The \emph{hot path}\index{hot path} of a trading system is the sequence of code executed between a triggering market
event and the resulting order: the code whose duration is the \omterm{def:ll:where-latency-comes-from:t2t}{tick-to-trade latency}. Everything else (logging to
disk, reporting, risk aggregation, configuration) runs off the hot path, on other threads or at other times.
\end{definition}

\Cref{fig:ll:where-latency-comes-from:path} is the map of Part~IV of this book, where each stage is built:
the feed handler that decodes and arbitrates the exchange's packets (chapter~18), the order-book builder
(chapter~19), the strategy engine (chapter~20), the pre-trade risk gate (chapter~22) and the order gateway
(chapter~21). Parts~I to~III supply what each stage is made of: the processor and its memory (chapters~2 to~5), the
languages (chapters~6 to~14) and the protocols (chapters~15 to~17).

Two physical quantities bound the path from below and are worth computing once.

\begin{example}[Light in fibre, bits on the wire]\label{ex:ll:where-latency-comes-from:physics}
A standard single-mode fibre has an effective group index of about 1.462 at a wavelength of 1\,550 nm (a
manufacturer's specification sheet): light in it covers a metre in $1.462/c = \qty{4.88}{\nano\second}$, so
the 30 metres of cable between a server and an exchange's switch inside one data centre cost about
\qty{146}{\nano\second} each way. Putting a frame on the wire takes its length in bits divided by the line rate,
plus 20 bytes of preamble, start delimiter and inter-frame gap: a minimal 64-byte frame occupies a 10~Gb/s link for
$84 \times 8 / 10 = \qty{67.2}{\nano\second}$ and a 25~Gb/s link for \qty{26.9}{\nano\second}. A software path of a
few microseconds is therefore dominated by computation and by the operating system, not by physics, until it is
very good.
\end{example}

\section{Latency is a distribution}

The \omterm{def:ll:where-latency-comes-from:t2t}{tick-to-trade latency} of one message is a draw. Its median describes the ordinary message; the messages that
cost money are often the others.

\begin{definition}[Tail latency, jitter]\label{def:ll:where-latency-comes-from:tail}
The \emph{tail latency}\index{tail latency} of a system is the upper part of its \omterm{def:ll:where-latency-comes-from:latency}{latency} distribution, reported as
high percentiles: the 99th (p99), the 99.9th (p99.9), the maximum over a period. \emph{Jitter}\index{jitter} is
the variability of \omterm{def:ll:where-latency-comes-from:latency}{latency} from one occurrence to the next, reported as a spread (p99 minus p50, or a standard
deviation); a system with low jitter is one whose tail is close to its median.
\end{definition}

We write $X$ for a \omterm{def:ll:where-latency-comes-from:latency}{latency}, $X_i$ for the \omterm{def:ll:where-latency-comes-from:latency}{latency} of stage $i$ of a path, and $q_p(X)$ for its $p$-quantile, so that
p99 is $q_{0.99}(X)$. The end-to-end \omterm{def:ll:where-latency-comes-from:latency}{latency} of a path is $X = \sum_i X_i$. Means add: $\E[X] = \sum_i \E[X_i]$.
Quantiles do not, and the direction of the error is not fixed.

\begin{proposition}[Tails of a sum]\label{prop:ll:where-latency-comes-from:union}
For any random variables $X_1,\dots,X_n$, whatever their dependence, and levels $\alpha_i > 0$ with $\sum_i \alpha_i < 1$,
\[
q_{1-\sum_i \alpha_i}\Big(\sum_i X_i\Big) \le \sum_i q_{1-\alpha_i}(X_i).
\]
In particular, stage targets $b_i$ met at the level $1-(1-p)/n$ guarantee an end-to-end target $\sum_i b_i$ at
level $p$.
\end{proposition}

\begin{proof}
Let $b_i = q_{1-\alpha_i}(X_i)$, so $\P(X_i > b_i) \le \alpha_i$. If $\sum_i X_i > \sum_i b_i$, then $X_i > b_i$ for
at least one $i$; hence $\P(\sum_i X_i > \sum_i b_i) \le \sum_i \P(X_i > b_i) \le \sum_i \alpha_i$, and the
$(1-\sum_i\alpha_i)$-quantile of the sum is at most $\sum_i b_i$. With $\alpha_i = (1-p)/n$ the level is $p$.
\end{proof}

The proposition says nothing about $q_{0.99}(\sum X_i)$ against $\sum q_{0.99}(X_i)$, and that comparison can go
either way. It fails in the direction that matters when each stage has rare stalls that none of the stage
percentiles sees.

\begin{example}[Two clean stages, one stalled path]\label{ex:ll:where-latency-comes-from:nonsub}
Two independent stages each take \qty{1}{\micro\second}, except that each stalls for \qty{100}{\micro\second} with
probability 0.9\%. Each stage's p99 is \qty{1}{\micro\second}: its stalls are rarer than 1\%. The path stalls when
either stage does, with probability $1-0.991^2 = 1.79\%$, so its p99 is \qty{101}{\micro\second}, fifty times the
sum of the stage p99s. (The same failure of additivity for value at risk is why One Quant Book~6 prefers expected
shortfall for aggregating risk.)
\end{example}

\Cref{fig:ll:where-latency-comes-from:ccdf} shows the same effect on a synthetic six-stage path, each stage with a
lognormal body and a 0.3\% chance of a stall of \qty{8}{\micro\second} on average. The stages' medians add up
almost exactly to the path's median (\qty{2.13}{\micro\second} against \qty{2.17}{\micro\second}) and the means add
exactly (\qty{2.32}{\micro\second}). The path is stalled with probability $1-0.997^6 \approx 1.8\%$, so its p99
(\qty{6.56}{\micro\second}) sits in the stalls, while the sum of the stage p99s (\qty{3.59}{\micro\second}) does
not see them. The union bound at level $1-0.01/6 \approx 99.83\%$ per stage gives \qty{28.7}{\micro\second}: safe,
and four times too pessimistic. Only the end-to-end distribution, measured message by message, gives the answer.

\begin{omfigure}
\begin{tikzpicture}
\begin{axis}[width=11.5cm,height=5.6cm,xmode=log,ymode=log,log origin y=infty,xmin=1.5,xmax=60,ymin=1e-4,ymax=1.2,
  xlabel={end-to-end latency $x$ (\si{\micro\second}, log scale)},ylabel={$\P(X > x)$ (log scale)},
  tick label style={font=\small},label style={font=\small},grid=major,grid style={gray!25},
  xtick={2,5,10,20,50},xticklabels={2,5,10,20,50},
  legend image code/.code={\draw[#1] (0cm,0cm) -- (0.6cm,0cm);},
  legend columns=2,legend style={font=\footnotesize,at={(0.5,-0.3)},anchor=north,draw=none,fill=none,
  /tikz/every even column/.append style={column sep=8pt}},table/col sep=comma]
\addplot[omDef,very thick] table[x=x_us,y=ccdf]{figdata/low-latency/01-where-latency-comes-from/ccdf.csv};
\addplot[black!45,dashed,thin,domain=1.5:60,samples=2] {0.01};
\addplot[omThm,very thick,ycomb] table[x=sum_p99,y=top]{figdata/low-latency/01-where-latency-comes-from/marks.csv};
\addplot[omMeth,very thick,ycomb] table[x=p99,y=top]{figdata/low-latency/01-where-latency-comes-from/marks.csv};
\addplot[omProp,very thick,dashed,ycomb] table[x=sum_union,y=top]{figdata/low-latency/01-where-latency-comes-from/marks.csv};
\legend{end-to-end tail,1\% level,sum of the stage p99s,end-to-end p99,{union bound (stages at 99.83\%)}}
\end{axis}
\end{tikzpicture}

\omcaption{The tail of a synthetic six-stage path whose stages each stall 0.3\% of the time. The sum of the stage
p99s misses the stalls and understates the path's p99; the union bound of
\cref{prop:ll:where-latency-comes-from:union} overstates it. Data: \texttt{ll\_budget.composition}, 400\,000 simulated
messages.}\label{fig:ll:where-latency-comes-from:ccdf}
\end{omfigure}

The same arithmetic applies to a real path, measured. \Cref{fig:ll:where-latency-comes-from:measured} times the four
stages of a toy path, Book~1's decoder and book with a trivial decision and an order encoder, message by message on
this book's laptop. The book, built on balanced trees and a hash map, is the largest stage at the median; the p99.9
of the decoder, the book and the decision is several times their p99, and the maximum is hundreds of times the
median: an interruption of the thread by the operating system, which this laptop, with no isolated cores, does not
prevent (chapter~13).

\begin{omfigure}
\begin{tikzpicture}
\begin{axis}[width=11.5cm,height=5.6cm,ybar,bar width=7pt,ymode=log,log origin y=infty,ymin=10,ymax=100000,
  symbolic x coords={decode,book,decide,encode,total},xtick=data,enlarge x limits=0.14,
  ylabel={nanoseconds (log scale)},tick label style={font=\small},label style={font=\small},
  grid=major,grid style={gray!25},
  legend columns=5,legend style={font=\footnotesize,at={(0.5,-0.2)},anchor=north,draw=none,fill=none,
  /tikz/every even column/.append style={column sep=6pt}},table/col sep=comma]
\addplot[fill=omDef!70,draw=omDef] table[x=stage,y=p50]{figdata/low-latency/01-where-latency-comes-from/measured_path.csv};
\addplot[fill=omDef!35,draw=omDef] table[x=stage,y=p99]{figdata/low-latency/01-where-latency-comes-from/measured_path.csv};
\addplot[fill=omThm!45,draw=omThm] table[x=stage,y=p999]{figdata/low-latency/01-where-latency-comes-from/measured_path.csv};
\addplot[fill=gray!40,draw=gray] table[x=stage,y=max]{figdata/low-latency/01-where-latency-comes-from/measured_path.csv};
\legend{p50,p99,p99.9,maximum}
\end{axis}
\end{tikzpicture}

\omcaption{Per-message \omterm{def:ll:where-latency-comes-from:latency}{latency} of the stages of a toy tick-to-trade path (decode one message, apply it to Book~1's
book, decide, encode an order), over 200 replays of Book~1's 2\,020-message sample, pinned to one CPU. Measured on a
laptop (Intel Core Ultra 7 155H) under WSL2, no isolated cores; each figure includes the timer's own cost of about
\qty{15}{\nano\second}. Data: \texttt{bench\_path.py}.}\label{fig:ll:where-latency-comes-from:measured}
\end{omfigure}

\section{Budgets from wire to wire}

\begin{definition}[Latency budget]\label{def:ll:where-latency-comes-from:budget}
A \emph{latency budget}\index{latency budget} is a target for the end-to-end \omterm{def:ll:where-latency-comes-from:latency}{latency} of a path at stated
percentiles, split into targets for each stage, against which every change to the system is measured.
\end{definition}

A budget turns ``make it faster'' into a list of numbers that a stage either meets or does not. It is written from
the outside in: what the strategy needs, what the physics and the hardware cost, and what is left for software.

\begin{method}[Writing a latency budget]\label{met:ll:where-latency-comes-from:budget}
\begin{enumerate}
\item State the requirement from the strategy's side: the \omterm{def:ll:where-latency-comes-from:latency}{latency} at which its edge disappears, and the
percentiles that matter (the median for races; the tail for exposure to stale quotes).
\item Subtract the fixed costs: cable and fibre lengths, serialisation, switches, the exchange's own gateway.
\item Split the remainder among the software stages in proportion to what they must do, at a stage level stricter
than the end-to-end level (\cref{prop:ll:where-latency-comes-from:union}).
\item Instrument every stage boundary (chapter~26) and check each release against the budget, stage by stage and end
to end, on the joint per-message samples (chapter~25).
\end{enumerate}
\end{method}

\begin{dated}{2026-09}{Published tick-to-trade figures}\label{dat:ll:where-latency-comes-from:published}
At CppCon 2017 an engineer of a large electronic market maker gave ``around \qty{2.5}{\micro\second}'' as a very
good minimum wire-to-wire time for a software-based trading system. The STAC-T0 benchmark measures the time from the
last bit of the inbound data needed for a decision to the first bit of the outbound order, with UDP in, TCP out and
essentially no trading logic; a report of June 2024 on a programmable-hardware card gave a minimum of
\qty{13.9}{\nano\second} for 507-byte frames. The two figures bracket the industry: software at microseconds,
hardware at tens of nanoseconds (One Quant Book~14, chapters~6 and~7).
\end{dated}

\begin{example}[A three-microsecond budget]\label{ex:ll:where-latency-comes-from:budget}
A software system must answer within \qty{3}{\micro\second} wire to wire at the median. The card and kernel-bypass
receive path take about \qty{0.9}{\micro\second} and the transmit path \qty{0.8}{\micro\second} (the synthetic
figures of \cref{fig:ll:where-latency-comes-from:ccdf}); \qty{1.3}{\micro\second} is left for decode, book, strategy,
risk and encode. Split \qty{60}{\nano\second}, \qty{120}{\nano\second}, \qty{200}{\nano\second},
\qty{50}{\nano\second} and the rest in reserve, it fits; the toy book of \cref{fig:ll:where-latency-comes-from:measured}
alone, at about \qty{80}{\nano\second} at the median but near a microsecond at p99.9, does not fit a tail
target. Chapter~19 builds one that does.
\end{example}

\section{Which strategies pay for which microseconds}

A race between two firms reacting to the same event is won by the one whose \omterm{def:ll:where-latency-comes-from:latency}{latency} on that message is lower. With
our \omterm{def:ll:where-latency-comes-from:latency}{latency} $X_A$ and the competitor's $X_B$, independent,
\[
\P(\text{win}) = \P(X_A < X_B),
\]
which depends on both distributions entirely, not on a single number. The median decides most races; the tail
decides the others.

\begin{example}[Median against tail in a race]\label{ex:ll:where-latency-comes-from:race}
Our \omterm{def:ll:where-latency-comes-from:latency}{latency} has a lognormal body with median \qty{2.0}{\micro\second} and a 2\% chance of a stall of
\qty{30}{\micro\second} on average; the competitor's median is \qty{2.2}{\micro\second} with stalls of the same
length 0.5\% of the time. Simulated over a million races, we win 66.1\%. Halving our median wins 98.0\%. Halving our
stall rate, which cuts our p99 from \qty{23.6}{\micro\second} to \qty{3.5}{\micro\second}, wins 66.7\%: a large
improvement of a percentile that races barely see (\cref{fig:ll:where-latency-comes-from:race}).
\end{example}

\begin{omfigure}
\begin{tikzpicture}
\begin{axis}[width=11cm,height=5.4cm,xmin=1.6,xmax=2.8,ymin=0,ymax=1,ytick={0.2,0.4,0.6,0.8,1},
  xlabel={our median latency (\si{\micro\second}); the competitor's is 2.2},ylabel={$\P(\text{win})$},
  tick label style={font=\small},label style={font=\small},grid=major,grid style={gray!25},
  legend columns=2,legend style={font=\footnotesize,at={(0.5,-0.3)},anchor=north,draw=none,fill=none,
  /tikz/every even column/.append style={column sep=8pt}},table/col sep=comma]
\addplot[omDef,very thick] table[x=median_us,y=with_stalls]{figdata/low-latency/01-where-latency-comes-from/race.csv};
\addplot[omThm,thick,dashed] table[x=median_us,y=no_stalls]{figdata/low-latency/01-where-latency-comes-from/race.csv};
\legend{with our 2\% stalls,without stalls}
\end{axis}
\end{tikzpicture}

\omcaption{Probability of winning a race against a competitor with median \qty{2.2}{\micro\second} and 0.5\% stalls,
as our median varies. Removing all our stalls moves the curve by about a point; moving our median by a tenth of a
microsecond moves it by about eight. Data: \texttt{ll\_budget.race\_curve}, 400\,000 simulated races per
point.}\label{fig:ll:where-latency-comes-from:race}
\end{omfigure}

The tail matters where the question is not who arrives first but how long a firm is exposed. A market maker's
quotes rest in the book while its view of the market is stale; for the duration of a stall they can be traded
against by anyone who has seen the new price. The expected cost of a stall is the rate of informed arrivals times
the stall's length times the loss per fill: adverse selection (One Quant Book~1, chapter~1), scaled by the tail. A
latency-arbitrage strategy (One Quant Book~11, chapter~9) pays for its median; a market maker pays for both, and
for its tail on the worst days, when stalls and news arrive together. A firm executing large orders over hours pays
for neither at the microsecond scale: its \omterm{def:ll:where-latency-comes-from:budget}{latency budget} is written in milliseconds, and its engineers' time is
better spent elsewhere.

\section{Tutorial: timing a path and composing its budget}\label{tut:ll:where-latency-comes-from:tut}

\begin{tutorial}
\textbf{Goal.} Time the stages of a toy tick-to-trade path message by message, then compose and check a budget with
\texttt{firm.latbudget}. \textbf{End state:} \cref{fig:ll:where-latency-comes-from:measured,fig:ll:where-latency-comes-from:ccdf}
and a budget report.
\begin{tutsteps}
\item \textbf{Time each stage.} The benchmark replays Book~1's sample and reads the processor's time-stamp counter
between stages (chapter~2 explains the counter and chapter~5 its calibration); the first pass warms the caches and is
discarded.
\omcode{code/low-latency/01-where-latency-comes-from/cpp/ll_path_bench.cpp}{28}{48}{Timing four stages per message with the time-stamp counter.}\label{lst:ll:where-latency-comes-from:bench}
\item \textbf{Run it pinned:} \texttt{python bench\_path.py} compiles the benchmark, runs it on one CPU with
\texttt{taskset}, and writes \texttt{measured\_path.csv} with a \texttt{.meta} file that records the machine, the
compiler, the flags and the load.
\item \textbf{Check a budget.} Stage targets and end-to-end targets are data; the check composes the stages jointly,
message by message.
\omcode{code/firm/latbudget/firm_latbudget.py}{44}{54}{A budget checked against aligned per-stage samples.}\label{lst:ll:where-latency-comes-from:check}
\item \textbf{Compose synthetic stages} with \texttt{ll\_budget.composition()} and compare the sum of stage percentiles
with the end-to-end percentile.
\end{tutsteps}
\textbf{What to change next.} Run the benchmark without \texttt{taskset} and next to a compilation of this book, and
compare the p99.9; replace the independent stalls of \texttt{ll\_budget} by stalls that hit every stage of a message at
once, and see which percentiles move.
\end{tutorial}

\section{Build: the latency budget}\label{bld:ll:where-latency-comes-from:latbudget}

\begin{build}
\bfield{Purpose} The contract every stage of the firm's \omterm{def:ll:where-latency-comes-from:hot}{hot path} is measured against, from this chapter to the
assembled path of chapter~26.
\bfield{Interface} \texttt{Stage(name, targets)}, \texttt{Budget(stages, end\_to\_end)},
\texttt{Budget.check(samples)} returning one row per stage and level (target, measured, ok),
\texttt{compose\_joint}, \texttt{compose\_independent(samples, n, seed)}, \texttt{stage\_level(p, n)},
\texttt{allocate(total, weights)}, \texttt{quantile(x, p)} (nearest rank).
\bfield{Rules} Samples are nanoseconds, one array per stage, aligned message by message when they come from one run;
end-to-end figures are composed jointly from aligned samples, never by adding percentiles; a stage target at level
$p$ for an end-to-end level $p$ is flagged as unsound.
\bfield{Acceptance tests} \texttt{code/firm/latbudget/tests/}: the union bound on independent, comonotone and
countermonotone pairs; the non-subadditive p99 of \cref{ex:ll:where-latency-comes-from:nonsub}; a budget with one
failing stage.
\bfield{Stretch} Confidence intervals on each measured percentile (bootstrap, One Quant Book~4, chapter~13); a
budget in cycles as well as nanoseconds.
\end{build}

\begin{omsources}
\item M.~Aquilina, E.~Budish and P.~O'Neill, ``Quantifying the high-frequency trading `arms race''', \emph{Quarterly
Journal of Economics} 137(1), 2022, 493--564.
\item C.~Cook, ``When a microsecond is an eternity: high performance trading systems in C++'', CppCon 2017.
\item STAC, STAC-T0 benchmark reports (tick-to-trade network I/O), 2020--2024.
\item Corning, SMF-28 ULL optical fibre, product information PI1470 (2020).
\end{omsources}

\section{Exercises}

\begin{exercise}[$\star$]\label{exo:ll:where-latency-comes-from:1}
A server sits 45 metres of fibre from the exchange's switch. What does the fibre cost each way, at a group index of
1.462?
\end{exercise}

\begin{exercise}[$\star$]\label{exo:ll:where-latency-comes-from:2}
How long does a 128-byte frame occupy a 10~Gb/s link, counting 20 bytes of preamble, delimiter and gap? And a 25~Gb/s
link?
\end{exercise}

\begin{exercise}[$\star$]\label{exo:ll:where-latency-comes-from:3}
Four stages have mean latencies of 300, 80, 150 and \qty{500}{\nano\second}. What is the mean end-to-end \omterm{def:ll:where-latency-comes-from:latency}{latency}, and
what can be said about its median?
\end{exercise}

\begin{exercise}[$\star\star$]\label{exo:ll:where-latency-comes-from:4}
A path has five independent stages, each stalling with probability 0.4\% per message. What is the probability that a
message meets at least one stall? Is the path's p99 in the stalls?
\end{exercise}

\begin{exercise}[$\star\star$]\label{exo:ll:where-latency-comes-from:5}
An end-to-end p99.9 target of \qty{10}{\micro\second} is to be guaranteed by targets on four stages. At what
percentile must each stage meet its target, and how must the targets relate?
\end{exercise}

\begin{exercise}[$\star\star$]\label{exo:ll:where-latency-comes-from:6}
Read \cref{fig:ll:where-latency-comes-from:race}. By how much must our median fall to gain as much as removing all
our stalls gains at a median of \qty{2.0}{\micro\second}?
\end{exercise}

\begin{exercise}[$\star\star\star$]\label{exo:ll:where-latency-comes-from:7}
\emph{Coding.} In \texttt{ll\_budget}, make the stalls common: when the path stalls (probability 1.8\% per message),
put the stall on one stage chosen at random, and none on the others. Which of the figures of
\cref{fig:ll:where-latency-comes-from:ccdf} change: the stage p99s, the end-to-end p99?
\end{exercise}

\begin{exercise}[$\star\star\star$]\label{exo:ll:where-latency-comes-from:8}
\emph{Find the flaw.} ``Our feed handler's p99 is \qty{400}{\nano\second}, our strategy's is \qty{600}{\nano\second}
and our gateway's is \qty{300}{\nano\second}, measured separately in three benchmarks: our tick-to-trade p99 is
\qty{1.3}{\micro\second}.''
\end{exercise}

\section{Problem: The Microsecond That Lost the Race}

\begin{problem}[{Weekend problem --- two firms, one stale quote}]\label{pb:ll:where-latency-comes-from:1}
Our firm and one competitor race for the same stale quotes. Our \omterm{def:ll:where-latency-comes-from:latency}{latency} has a lognormal body (median
\qty{2.0}{\micro\second}, dispersion 0.15) and a 2\% chance per message of a stall of exponential length with mean
\qty{30}{\micro\second}; the competitor's median is \qty{2.2}{\micro\second} and it stalls 0.5\% of the time.

\medskip\noindent\textbf{Part I --- The race.}
\begin{enumerate}
\item Write the probability that we win a race in terms of the two \omterm{def:ll:where-latency-comes-from:latency}{latency} distributions. What does it assume about
their dependence?
\item Without stalls, both bodies lognormal with the same dispersion, what is $\P(\text{win})$ in closed form?
Evaluate it.
\item What share of races do we win with the stalls (the simulation)?
\item Why is the answer below the no-stall answer by less than our stall rate?
\end{enumerate}

\medskip\noindent\textbf{Part II --- Median or tail.}
\begin{enumerate}[resume]
\item What are our p50 and p99 with the stalls?
\item What does halving our median win? What does halving our stall rate win, and what does it do to our p99?
\item Removing all stalls wins how much?
\item Which should the firm pay for, if it only races?
\end{enumerate}

\medskip\noindent\textbf{Part III --- Exposure.}
\begin{enumerate}[resume]
\item Our quotes are exposed while we are stalled. With stalls on 2\% of messages and a mean of
\qty{30}{\micro\second}, what fraction of time is our view of the market stale if a message arrives every
\qty{10}{\micro\second} on average and stalls do not overlap?
\item If informed takers arrive at 20 a second and each fill on a stale quote costs half a tick of \$0.01 on 100
shares, what does the staleness cost a day of 6.5 hours?
\item How does halving the stall rate change that cost?
\item Which part of the distribution does the exposure depend on, and which the race?
\end{enumerate}

\medskip\noindent\textbf{Part IV --- The budget.}
\begin{enumerate}[resume]
\item The firm wants $\P(\text{win}) \ge 0.9$ against this competitor. Read the median it needs from
\cref{fig:ll:where-latency-comes-from:race}.
\item Write that as a budget over five software stages with equal shares of what remains after
\qty{1.0}{\micro\second} of network and kernel. What does each stage get?
\item At what level must each stage meet its target to guarantee an end-to-end p99?
\item State the \emph{named result}: our win probability, and its gain from halving the median against halving the
stall rate.
\item Why does a single number (``our \omterm{def:ll:where-latency-comes-from:latency}{latency} is \qty{2}{\micro\second}'') describe neither race nor exposure?
\item What measurement would tell the firm the competitor's distribution?
\item Name one strategy for which neither the median nor the tail at microseconds matters.
\item In one sentence: what does a \omterm{def:ll:where-latency-comes-from:budget}{latency budget} buy an engineering team?
\end{enumerate}
\end{problem}

\section{Interview questions}

\begin{interviewq}[$\star$ \iqroles{developer}]\label{iq:ll:where-latency-comes-from:1}
What is the difference between tick-to-trade and \omterm{def:ll:where-latency-comes-from:t2t}{wire-to-wire latency}, and how would you measure each?
\end{interviewq}

\begin{interviewq}[$\star$ \iqroles{developer, trader}]\label{iq:ll:where-latency-comes-from:2}
Why do trading firms report p99 and p99.9 rather than the mean \omterm{def:ll:where-latency-comes-from:latency}{latency}?
\end{interviewq}

\begin{interviewq}[$\star\star$ \iqroles{developer, researcher}]\label{iq:ll:where-latency-comes-from:3}
Three components have p99 latencies of 1, 2 and \qty{3}{\micro\second}. What can you say about the p99 of the
pipeline?
\end{interviewq}

\begin{interviewq}[$\star\star$ \iqroles{developer}]\label{iq:ll:where-latency-comes-from:4}
Light travels at about \qty{5}{\nano\second} a metre in fibre. Where else does time go between a packet reaching a
server and an order leaving it?
\end{interviewq}

\begin{interviewq}[$\star\star$ \iqroles{trader, researcher}]\label{iq:ll:where-latency-comes-from:5}
Your firm wins 60\% of races against its main competitor. Should it spend its next quarter on the median or on the
tail of its \omterm{def:ll:where-latency-comes-from:latency}{latency}?
\end{interviewq}

\begin{interviewq}[$\star\star\star$ \iqroles{developer}]\label{iq:ll:where-latency-comes-from:6}
Design the measurement that would let you write a \omterm{def:ll:where-latency-comes-from:budget}{latency budget} for an existing trading system you did not build.
\end{interviewq}
